\documentclass[11pt,a4paper]{article}

\usepackage[margin=2.5cm]{geometry}
\usepackage[T1]{fontenc}
\usepackage[utf8]{inputenc}
\usepackage{lmodern}
\usepackage{amsmath,amssymb,amsthm,mathtools}
\usepackage{enumitem}
\usepackage{booktabs,tabularx,array}
\usepackage{parskip}
\usepackage{url}
\usepackage{xcolor}
\usepackage[colorlinks=true,linkcolor=blue!50!black,urlcolor=blue!50!black,citecolor=blue!50!black]{hyperref}

% ---------------------------------------------------------------- numbering
\numberwithin{equation}{section}
\setcounter{section}{-1}

% ------------------------------------------- theorem environments (W.1, ...)
\theoremstyle{plain}
\newtheorem{theorem}{Theorem}
\renewcommand{\thetheorem}{W.\arabic{theorem}}
\newtheorem{corollary}[theorem]{Corollary}
\newtheorem{proposition}[theorem]{Proposition}
\theoremstyle{remark}
\newtheorem{remark}[theorem]{Remark}

% ------------------------------------------------------------------- macros
\newcommand{\opnorm}[1]{{\vert\kern-0.25ex\vert\kern-0.25ex\vert #1 \vert\kern-0.25ex\vert\kern-0.25ex\vert}}
\DeclareMathOperator{\vect}{vec}
\DeclareMathOperator{\sign}{sign}
\DeclareMathOperator{\diag}{diag}
\DeclareMathOperator{\tr}{tr}
\DeclareMathOperator*{\argmin}{arg\,min}
\newcommand{\changed}{$\blacktriangleright$\ }
\DeclareUrlCommand\code{\urlstyle{tt}}

\title{\textbf{Adaptive (weighted) lasso in GCLMs}\\[0.4em]}
\author{}
\date{}

\begin{document}
\maketitle

% =========================================================================
\section{Setting (DDK \S\S2--3)}\label{sec:setting}

\paragraph{Data and model.} $X_1,\dots,X_n\in\mathbb R^p$ are i.i.d.\ $N(0,\Sigma^*)$ with $M^*\Sigma^*+\Sigma^*M^{*\top}+C=0$, where $M^*\in\mathrm{Stab}_p$ and $C\in PD_p$ is known. $\hat\Sigma=\frac1n\sum_iX_iX_i^\top$.

\paragraph{Design and Gram matrix.} For symmetric $\Sigma$:
\[
A(\Sigma)=(\Sigma\otimes I_p)+(I_p\otimes\Sigma)K^{(p,p)},\qquad
\Gamma(\Sigma)=A(\Sigma)^\top A(\Sigma),\qquad
g(\Sigma)=-A(\Sigma)^\top\vect(C),
\]
so that
\[
\tfrac12\|M\Sigma+\Sigma M^\top+C\|_F^2=\tfrac12\vect(M)^\top\Gamma(\Sigma)\vect(M)-g(\Sigma)^\top\vect(M)+\tfrac12\|C\|_F^2 .
\]

\paragraph{Shorthand.} $\hat\Gamma=\Gamma(\hat\Sigma)$, $\Gamma^*=\Gamma(\Sigma^*)$, $\hat g$, $g^*$, $\Delta_\Gamma=\hat\Gamma-\Gamma^*$, $\Delta_g=\hat g-g^*$. Indices are pairs $(j,k)$; $m=\vect(M)$, $m^*=\vect(M^*)$.

\paragraph{Support.} $S=\{(j,k):M^*_{jk}\ne0\}$, $d=|S|$; Since $M^*$ solves the Lyapunov equation, $\Gamma^*m^*=g^*$.

\paragraph{Constants (DDK).}
\begin{gather*}
c_{\Gamma^*}=\opnorm{(\Gamma^*_{SS})^{-1}}_\infty,\qquad c_{M^*}=\|m^*\|_\infty,\\
P:=\Gamma^*_{S^cS}(\Gamma^*_{SS})^{-1},\qquad \rho:=\opnorm{P}_\infty,\qquad \beta_{\min}:=\min_{(j,k)\in S}|M^*_{jk}|
\end{gather*}


% =========================================================================
\section{Estimator, weights, weighted irrepresentability constant}\label{sec:estimator}

\paragraph{Pilot and weights}
\begin{itemize}
  \item The random pilot $\hat M^0$ is the lasso (diagonals penalized), at the smallest $\lambda$ of its grid.
  \item With $\hat\mu:=\max_{(k,l)}|\hat M^0_{kl}|$, the maximum over all entries including the diagonal,
  \[
  \hat w_{jk}=\frac{\hat\mu}{|\hat M^0_{jk}|}\quad\text{if }\hat M^0_{jk}\ne0,\qquad
  \hat w_{jk}=\infty\ \ (\text{entry excluded, stays }0)\quad\text{if }\hat M^0_{jk}=0 .
  \]
  \item All finite weights are $\ge1$, with equality at the largest pilot entry.
\end{itemize}

\paragraph{Estimator} (weighted direct Lyapunov lasso, all entries penalised):
\begin{equation}\label{eq:est}
\begin{split}
\hat M&\in\argmin_{M\in\mathbb R^{p\times p}}\ \tfrac12\|M\hat\Sigma+\hat\Sigma M^\top+C\|_F^2+\lambda\sum_{j,k}\hat w_{jk}|M_{jk}|\\
&=\argmin_m\ \tfrac12m^\top\hat\Gamma m-\hat g^\top m+\lambda\|\hat Wm\|_1,\qquad \hat W=\diag(\hat w).
\end{split}
\end{equation}
An infinite weight means the constraint $m_j=0$.

\paragraph{Scale invariance.} Replacing $(\lambda,\hat w)$ by $(\lambda/c,\,c\hat w)$ leaves the objective unchanged. The products $\lambda\hat w_j$ matter.

\paragraph{Quantities.} Define
\[
\underline w:=\min_{j\in S^c}\hat w_j,\qquad
\bar w_S:=\max_{k\in S}\hat w_k<\infty,\qquad
r:=\frac{\bar w_S}{\underline w}=\frac{\max_{j\in S^c}|\hat M^0_j|}{\min_{k\in S}|\hat M^0_k|},
\]

\begin{itemize}
  \item $\bar w_S<\infty$ requires the pilot to have no zero on $S$, a necessary condition for recovery. In particular, all infinite weights lie in $S^c$.
 \end{itemize}

\paragraph{Weighted irrepresentability constant} (Zhou, van de Geer \& Bühlmann 2009, Def.~8.1):
\begin{equation}\label{eq:rhow}
\rho_w:=\max_{j\in S^c}\ \frac{1}{\hat w_j}\sum_{k\in S}|P_{jk}|\,\hat w_k\ \ \Big(=\opnorm{\hat W_{S^c}^{-1}P_{S^c\cdot}\hat W_{SS}}_\infty\Big).
\end{equation}

\paragraph{A free constant.} Fix any $\kappa\ge1+\rho$. The natural choice is $\kappa=1+\rho$. In DDK's setting, $\rho\le1-\alpha$ and $\kappa=2-\alpha$ is admissible, and this is the choice that reproduces D.1.

% =========================================================================
\section{W.1 (deterministic; analogue of DDK Theorem D.1)}\label{sec:W1}

\begin{theorem}\label{thm:W1}
Let $M^*\in\mathrm{Stab}_p$ have support $S$, let $\hat w$ be weights with $\bar w_S<\infty$, and let $\kappa\ge1+\rho$. Assume:
\begin{enumerate}[label=\upshape(W\arabic*),leftmargin=3.2em]
  \item $\Gamma^*_{SS}$ is invertible;
  \item (weighted irrepresentability) $\rho_w\le1-\alpha$ for some $\alpha\in(0,1]$;
  \item $\opnorm{(\Delta_\Gamma)_{\cdot S}}_\infty\le\varepsilon_1$ and $\|\Delta_g\|_\infty\le\varepsilon_2$, with
  \[
  \varepsilon_1\le\frac{\alpha}{3c_{\Gamma^*}(\kappa r+\alpha)};
  \]
  \item the effective penalty off the support satisfies
  \[
  \lambda\,\underline w>\frac{3\kappa}{\alpha}\max\{c_{M^*}\varepsilon_1,\ \varepsilon_2\}.
  \]
\end{enumerate}
Then:
\begin{enumerate}[label=\upshape\textbf{(\alph*)},leftmargin=2.5em]
  \item The minimiser $\hat M$ of \eqref{eq:est} is unique, its support is contained in the true support ($\hat S\subseteq S$), and
  \begin{equation}\label{eq:bw}
  \begin{split}
  \|\hat M-M^*\|_\infty&\le b_w:=\frac{c_{\Gamma^*}}{1-c_{\Gamma^*}\varepsilon_1}\big(c_{M^*}\varepsilon_1+\varepsilon_2+\lambda\bar w_S\big)\\
  &<c_{\Gamma^*}\,\lambda\Big(\bar w_S+\frac{\alpha}{\kappa}\,\underline w\Big).
  \end{split}
  \end{equation}
  The last bound is informative only when $\underline w<\infty$.
  \item If moreover $\beta_{\min}=\min_{(j,k)\in S}|M^*_{jk}|>b_w$, then $\hat S=S$ and $\sign(\hat M_{jk})=\sign(M^*_{jk})$ for all $(j,k)\in S$.
\end{enumerate}
\end{theorem}

\paragraph{Reduction to DDK Theorem D.1.} Take $\hat w\equiv1$ and $\kappa=2-\alpha$; this needs $\rho\le1-\alpha$. Then:
\begin{itemize}
  \item (W2) is DDK's (D.1), $\rho\le1-\alpha$;
  \item (W3) is DDK's $\varepsilon_1\le\alpha/(6c_{\Gamma^*})$, since $3c_{\Gamma^*}(\kappa\cdot1+\alpha)=6c_{\Gamma^*}$;
  \item (W4) is DDK's $\lambda>\frac{3(2-\alpha)}{\alpha}\max\{c_{M^*}\varepsilon_1,\varepsilon_2\}$;
  \item the right-hand bound in \eqref{eq:bw} is $c_{\Gamma^*}\lambda\big(1+\frac{\alpha}{2-\alpha}\big)=\frac{2c_{\Gamma^*}\lambda}{2-\alpha}$, which is DDK's bound;
  \item DDK's beta-min $\beta_{\min}>2c_{\Gamma^*}\lambda/(2-\alpha)$ implies the one in (b).
\end{itemize}

% =========================================================================
\section{Proof of W.1}\label{sec:proof}

\paragraph{KKT conditions} Problem \eqref{eq:est} is convex, since $\hat\Gamma\succeq0$.
\begin{equation}\label{eq:kkt}
\begin{gathered}
\big(\hat\Gamma\hat m-\hat g\big)_j+\lambda\,\hat w_j\hat z_j=0\quad(\text{for every $j$ with $\hat{w}_j < \infty$}),\\
\hat z_j=\sign(\hat m_j)\ \text{if }\hat m_j\ne0,\qquad \hat z_j\in[-1,1]\ \text{if }\hat m_j=0 .
\end{gathered}
\end{equation}

\emph{Block form.} Since $\Gamma^*m^*=g^*$, we have $\hat\Gamma\hat m-\hat g=\Gamma^*(\hat m-m^*)+\Delta_\Gamma\hat m-\Delta_g$. With $\Delta_M:=\hat m-m^*$, the KKT conditions \eqref{eq:kkt} read
\begin{align*}
&\begin{bmatrix}\Gamma^*_{SS}&\Gamma^*_{SS^c}\\\Gamma^*_{S^cS}&\Gamma^*_{S^cS^c}\end{bmatrix}
\begin{bmatrix}(\Delta_M)_S\\(\Delta_M)_{S^c}\end{bmatrix}\\
&\quad+\begin{bmatrix}(\Delta_\Gamma)_{SS}&(\Delta_\Gamma)_{SS^c}\\(\Delta_\Gamma)_{S^cS}&(\Delta_\Gamma)_{S^cS^c}\end{bmatrix}
\begin{bmatrix}\hat m_S\\\hat m_{S^c}\end{bmatrix}
-\begin{bmatrix}(\Delta_g)_S\\(\Delta_g)_{S^c}\end{bmatrix}
+\lambda\begin{bmatrix}(\hat w\circ\hat z)_S\\(\hat w\circ\hat z)_{S^c}\end{bmatrix}
=\begin{bmatrix}0\\0\end{bmatrix},
\end{align*}
where the $S^c$-rows are taken only for $j$ with $\hat w_j<\infty$.

\changed The PDW constructs a pair $(\tilde m,\tilde z)$ satisfying \eqref{eq:kkt} with $\tilde m_{S^c}=0$.

\paragraph{Step 1 (restricted problem)} Let
\begin{equation}\label{eq:restricted}
\tilde m=\argmin_{m:\ m_{S^c}=0}\ \tfrac12m^\top\hat\Gamma m-\hat g^\top m+\lambda\|\hat Wm\|_1 .
\end{equation}

The restricted objective is strictly convex (as $\hat\Gamma$ positive semidefinite and invertible according to DDK) and $\tilde m$ is unique.


At $\tilde m$, where $\tilde m_{S^c}=0$, the $S$-rows give
\begin{equation}\label{eq:Srows}
(\hat\Gamma \tilde m - \hat g)_S + \lambda(\hat w\circ\tilde z)_S
  = \hat\Gamma_{SS}\tilde m_S - \hat g_S + \lambda(\hat w\circ\tilde z)_S = 0.
\end{equation}

\emph{\textbf{Error bound:}} 

Since $m^*_{S^c}=0$ and $\Gamma^*_{SS}m^*_S-g^*_S = 0$, we have
\begin{equation*}
\hat\Gamma_{SS}m^*_S-\hat g_S= \underbrace{\Gamma_{SS}^*m^*_S - g^*_S}_{=0} + (\Delta_\Gamma)_{SS}m^*_S-(\Delta_g)_S = (\Delta_\Gamma)_{SS}m^*_S-(\Delta_g)_S .
\end{equation*}
Subtracting this from \eqref{eq:Srows} gives
\begin{equation*}
\hat\Gamma_{SS}(\underbrace{\tilde{m}_S-m^*_S}_{=\tilde\Delta_S}) + \lambda(\hat w\circ\tilde z)_S = -((\Delta_\Gamma)_{SS}m^*_S - (\Delta_g)_S)
\end{equation*}
As $\hat\Gamma_{SS}$ invertible (shown in DDK):
\[
\tilde\Delta_S:=\tilde m_S-m^*_S=-(\hat\Gamma_{SS})^{-1}\underbrace{\Big((\Delta_\Gamma)_{SS}m^*_S-(\Delta_g)_S+\lambda(\hat w\circ\tilde z)_S\Big)}_{=:v}.
\]
Then: $\|\tilde\Delta_S\|_\infty\le\opnorm{(\hat\Gamma_{SS})^{-1}}_\infty\|v\|_\infty$. 

For $v$,
\begin{align*}
\|v\|_\infty
&\le\|(\Delta_\Gamma)_{SS}m^*_S\|_\infty+\|(\Delta_g)_S\|_\infty+\lambda\|(\hat w\circ\tilde z)_S\|_\infty && \text{(triangle inequality)}\\
&\le\opnorm{(\Delta_\Gamma)_{\cdot S}}_\infty\,c_{M^*}+\|\Delta_g\|_\infty+\lambda\bar w_S && (\|m^*_S\|_\infty\le c_{M^*},\ |\tilde z_k|\le1)\\
&\le c_{M^*}\varepsilon_1+\varepsilon_2+\lambda\bar w_S && \text{(W3)}.
\end{align*}

For the inverse, write $\hat\Gamma_{SS}=\Gamma^*_{SS}+(\Delta_\Gamma)_{SS}=\Gamma^*_{SS}(I+X)$ with $X:=(\Gamma^*_{SS})^{-1}(\Delta_\Gamma)_{SS}$. Then
\begin{align*}
\opnorm{X}_\infty
&\le\opnorm{(\Gamma^*_{SS})^{-1}}_\infty\,\opnorm{(\Delta_\Gamma)_{SS}}_\infty && \text{(submultiplicativity)}\\
&\le c_{\Gamma^*}\varepsilon_1 && \text{(W3)}\\
&\le\frac{\alpha}{3(\kappa r+\alpha)}\le\frac13 && \text{(W3)}.
\end{align*}
So the Neumann series $(I+X)^{-1}=\sum_{k\ge0}(-X)^k$ converges, $\hat\Gamma_{SS}$ is invertible, and
\begin{align}
\opnorm{(\hat\Gamma_{SS})^{-1}}_\infty
&=\opnorm{(I+X)^{-1}(\Gamma^*_{SS})^{-1}}_\infty\notag\\
&\le\opnorm{(I+X)^{-1}}_\infty\,c_{\Gamma^*} && \text{(submultiplicativity)}\notag\\
&\le\sum_{k\ge0}\opnorm{X}_\infty^k\,c_{\Gamma^*}=\frac{c_{\Gamma^*}}{1-\opnorm{X}_\infty} && \text{(Neumann series)}\notag\\
&\le\frac{c_{\Gamma^*}}{1-c_{\Gamma^*}\varepsilon_1} && (\opnorm{X}_\infty\le c_{\Gamma^*}\varepsilon_1).\label{eq:inv}
\end{align}
Multiplying the two bounds gives
\begin{equation}\label{eq:errbound}
\|\tilde\Delta_S\|_\infty\le\frac{c_{\Gamma^*}}{1-c_{\Gamma^*}\varepsilon_1}\big(c_{M^*}\varepsilon_1+\varepsilon_2+\lambda\bar w_S\big)=:b_w.
\end{equation}

\emph{\textbf{Converting the bound into $\lambda$:}}

By (W4), $c_{M^*}\varepsilon_1+\varepsilon_2\le2\max\{c_{M^*}\varepsilon_1,\varepsilon_2\}<\frac{2\alpha}{3\kappa}\lambda\underline w$. By (W3), $\frac{1}{1-c_{\Gamma^*}\varepsilon_1}\le\frac{3(\kappa r+\alpha)}{3\kappa r+2\alpha}$. Hence
\begin{equation}\label{eq:bwlambda}
\begin{split}
b_w&<c_{\Gamma^*}\cdot\frac{3(\kappa r+\alpha)}{3\kappa r+2\alpha}\cdot\lambda\underline w\Big(r+\frac{2\alpha}{3\kappa}\Big)\\
&=c_{\Gamma^*}\lambda\underline w\,\frac{\kappa r+\alpha}{\kappa}=c_{\Gamma^*}\lambda\Big(\bar w_S+\frac{\alpha}{\kappa}\underline w\Big).
\end{split}
\end{equation}

\changed This is (D.4). With $\hat w\equiv1$ and $\kappa=2-\alpha$ the right-hand side is DDK's $\frac{6c_{\Gamma^*}}{6-\alpha}\cdot\frac{6-\alpha}{3(2-\alpha)}\lambda=\frac{2c_{\Gamma^*}\lambda}{2-\alpha}$.

\paragraph{Step 2 (subgradient on $S$).} Take $\tilde z_S$ from \eqref{eq:Srows}. By construction it is a subgradient of $\|\cdot\|_1$ at $\tilde m_S$.

\paragraph{Step 3 (dual vector on $S^c$).} Evaluate the block form at $(\tilde m,\tilde z)$. Since $\tilde m_{S^c}=m^*_{S^c}=0$, we have $(\Delta_M)_{S^c}=0$ and $(\Delta_M)_S=\tilde\Delta_S$:
\begin{align*}
&\begin{bmatrix}\Gamma^*_{SS}&\Gamma^*_{SS^c}\\\Gamma^*_{S^cS}&\Gamma^*_{S^cS^c}\end{bmatrix}
\begin{bmatrix}\tilde\Delta_S\\0\end{bmatrix}\\
&\quad+\begin{bmatrix}(\Delta_\Gamma)_{SS}&(\Delta_\Gamma)_{SS^c}\\(\Delta_\Gamma)_{S^cS}&(\Delta_\Gamma)_{S^cS^c}\end{bmatrix}
\begin{bmatrix}\tilde m_S\\0\end{bmatrix}
-\begin{bmatrix}(\Delta_g)_S\\(\Delta_g)_{S^c}\end{bmatrix}
+\lambda\begin{bmatrix}(\hat w\circ\tilde z)_S\\(\hat w\circ\tilde z)_{S^c}\end{bmatrix}
=\begin{bmatrix}0\\0\end{bmatrix}.
\end{align*}
The zeros remove the $S^c$-columns of $\Gamma^*$ and $\Delta_\Gamma$. Collect the noise in
\[
u:=(\Delta_\Gamma)_{\cdot S}\tilde m_S-\Delta_g,\qquad\text{i.e.}\quad u_S=(\Delta_\Gamma)_{SS}\tilde m_S-(\Delta_g)_S,\quad u_{S^c}=(\Delta_\Gamma)_{S^cS}\tilde m_S-(\Delta_g)_{S^c}.
\]
The two block rows read
\begin{align*}
\Gamma^*_{SS}\tilde\Delta_S+u_S+\lambda(\hat w\circ\tilde z)_S&=0 && (S\text{-rows}),\\
\Gamma^*_{S^cS}\tilde\Delta_S+u_{S^c}+\lambda(\hat w\circ\tilde z)_{S^c}&=0 && (S^c\text{-rows with }\hat w_j<\infty).
\end{align*}
Solving the $S$-rows gives $\tilde\Delta_S=-(\Gamma^*_{SS})^{-1}\big(u_S+\lambda(\hat w\circ\tilde z)_S\big)$. 

Define $\tilde z_j$, for $j\in S^c$ with $\hat w_j<\infty$, by the $S^c$-rows, solved for the $\tilde z$-term. Then insert $\tilde\Delta_S$ and write $P=\Gamma^*_{S^cS}(\Gamma^*_{SS})^{-1}$:
\begin{align*}
\lambda(\hat w\circ\tilde z)_{S^c}
&=-\Gamma^*_{S^cS}\tilde\Delta_S-u_{S^c} && (S^c\text{-rows})\\
&=\Gamma^*_{S^cS}(\Gamma^*_{SS})^{-1}\big(u_S+\lambda(\hat w\circ\tilde z)_S\big)-u_{S^c} && (\text{insert }\tilde\Delta_S)\\
&=P\big(u_S+\lambda(\hat w\circ\tilde z)_S\big)-u_{S^c} && (\text{definition of }P)\\
&=Pu_S-u_{S^c}+\lambda P(\hat w\circ\tilde z)_S && (\text{expand}),
\end{align*}
and dividing row $j$ by $\lambda\hat w_j$,
\begin{equation}\label{eq:dual}
\tilde z_j=\underbrace{\frac{1}{\lambda\hat w_j}\big[Pu_S-u_{S^c}\big]_j}_{\text{noise}}+\underbrace{\frac{1}{\hat w_j}\big[P(\hat w\circ\tilde z)_S\big]_j}_{\text{deterministic}} .
\end{equation}

\changed This is (D.5) with $\sign(\vect(\tilde M)_S)$ replaced by $(\hat w\circ\tilde z)_S$ and row $j$ divided by $\hat w_j$.

\emph{Deterministic term.} Since $|\tilde z_k|\le1$,
\begin{equation}\label{eq:det}
\Big|\frac{1}{\hat w_j}\big[P(\hat w\circ\tilde z)_S\big]_j\Big|\le\frac{1}{\hat w_j}\sum_{k\in S}|P_{jk}|\hat w_k\le\rho_w\le1-\alpha\qquad\text{by (W2)}.
\end{equation}

\changed DDK bound this term by $\rho\le1-\alpha$. This is the only place where the weighted condition enters.

\emph{Noise term.}
\begin{align}
|\text{noise}_j|
&\le\frac{1}{\lambda\hat w_j}\big(\rho\|u_S\|_\infty+\|u_{S^c}\|_\infty\big) && (\opnorm{P}_\infty=\rho)\notag\\
&\le\frac{\kappa}{\lambda\underline w}\|u\|_\infty && (1+\rho\le\kappa,\ \hat w_j\ge\underline w)\notag\\
&\le\frac{\kappa}{\lambda\underline w}\big(\varepsilon_1\|\tilde m_S\|_\infty+\varepsilon_2\big) && \text{Def. of $u$, (W3)}\notag\\
&\le\underbrace{\frac{\kappa\,c_{M^*}\varepsilon_1}{\lambda\underline w}}_{G_1}+\underbrace{\frac{\kappa\,\varepsilon_1\|\tilde\Delta_S\|_\infty}{\lambda\underline w}}_{G_2}+\underbrace{\frac{\kappa\,\varepsilon_2}{\lambda\underline w}}_{G_3} && (\|\tilde m_S\|_\infty\le c_{M^*}+\|\tilde\Delta_S\|_\infty).\label{eq:noise}
\end{align}

\changed These are DDK's $G_1,G_2,G_3$ (pp.~532--533), with $(2-\alpha)/\lambda$ replaced by $\kappa/(\lambda\underline w)$. Each is below $\alpha/3$:
\begin{itemize}
  \item $G_1,G_3<\alpha/3$ by (W4);
  \item by \eqref{eq:errbound}--\eqref{eq:bwlambda}, $G_2<\frac{\kappa\varepsilon_1}{\lambda\underline w}\,c_{\Gamma^*}\lambda\underline w\,\frac{\kappa r+\alpha}{\kappa}=c_{\Gamma^*}\varepsilon_1(\kappa r+\alpha)\le\alpha/3$ by (W3).
\end{itemize}
With \eqref{eq:det},
\begin{equation}\label{eq:dualfeas}
\max_{j\in S^c,\ \hat w_j<\infty}|\tilde z_j|<(1-\alpha)+\tfrac\alpha3+\tfrac\alpha3+\tfrac\alpha3=1 .
\end{equation}
If $\hat w_j=\infty$ for every $j\in S^c$ (the pilot is zero on all of $S^c$), Step 3 is empty: every $S^c$-coordinate is fixed at $0$ by the constraint, so there is no $\tilde z_j$ to bound and $\hat m_{S^c}=0$ holds automatically.

\paragraph{Step 4 (the witness is the unique minimiser; $\hat S\subseteq S$).} The pair $(\tilde m,\tilde z)$ satisfies \eqref{eq:kkt}:
\begin{itemize}
  \item the $S$-rows by \eqref{eq:Srows} and the rows $j\in S^c$ with $\hat w_j<\infty$ by construction;
  \item $\tilde m_j=0$ wherever $\hat w_j=\infty$, since all infinite weights lie in $S^c$;
  \item $\tilde z$ is a valid subgradient, since $|\tilde z_j|<1$ at the zero coordinates $j\in S^c$ with $\hat w_j<\infty$ by \eqref{eq:dualfeas}, and $\tilde z_S$ comes from Step~2.
\end{itemize}
So $\tilde m$ is a minimiser.

\emph{\textbf{Every minimiser has support in $S$}} (Wainwright 2009, Lemma~1, cited by DDK on p.~532). Strict dual feasibility implies hat $\hat{S} \subseteq S$.

\emph{\textbf{Uniqueness.}} Every minimiser lives on $\{m_{S^c}=0\}$, where the objective is strictly convex (Step 1). Part (a) follows from \eqref{eq:errbound}--\eqref{eq:bwlambda}. \hfill$\square_{\text{(a)}}$

\paragraph{Part (b).} For $(j,k)\in S$, (a) gives $|\hat M_{jk}-M^*_{jk}|\le b_w<\beta_{\min}\le|M^*_{jk}|$. Hence $\hat M_{jk}\ne0$ with the sign of $M^*_{jk}$, so $S\subseteq\hat S$, and with (a), $\hat S=S$. 

% =========================================================================
\section{Theorem W.2 (probabilistic; the analogue of DDK Theorem 2)}\label{sec:W2}

\paragraph{DDK's tools, used unchanged.}
\begin{itemize}
  \item Lemma~E.6: $\opnorm{\Delta_\Sigma}_2<\min\{\varepsilon_1/(\sqrt d(4+8c_{\Sigma^*})),\ \varepsilon_2/(2c_C)\}$ implies $\opnorm{(\Delta_\Gamma)_{\cdot S}}_\infty<\varepsilon_1$ and $\|\Delta_g\|_\infty<\varepsilon_2$, with $c_{\Sigma^*}=\opnorm{\Sigma^*}_2$ and $c_C=\|\vect(C)\|_2$.
  \item Corollary~E.5 (from Wainwright 2019, Thm~6.5) bounds $\opnorm{\Delta_\Sigma}_2$.
  \item Neither involves the penalty. As in DDK,
  \[
  \begin{split}
  \tilde c=\max\Big\{&\frac{4\max\{1,c_{\Sigma^*}^2\}(4+8c_{\Sigma^*})^2}{c_3},\ 16c_1^2c_{\Sigma^*}^2(4+8c_{\Sigma^*})^2,\\
  &\frac{16\max\{1,c_{\Sigma^*}^2\}c_C^2}{c_3},\ 64c_1^2c_{\Sigma^*}^2c_C^2\Big\},
  \end{split}
  \]
  with the universal constants $c_1,c_2,c_3$ of Theorem~E.4.
\end{itemize}

\begin{theorem}\label{thm:W2}
In the setting of DDK Theorem~2, let $\kappa\ge1+\rho$, $\alpha\in(0,1]$ and $\bar r\ge0$. Let $\mathcal E_w$ be an event on which the weights satisfy $\bar w_S<\infty$, $\rho_w\le1-\alpha$ and $r\le\bar r$. Put
\[
\varepsilon:=\sqrt{\tau_1\tilde c\,d\,p/n},\qquad c_w:=\frac{3c_{\Gamma^*}(\kappa\bar r+\alpha)}{\alpha}\qquad\Big(\text{DDK: }c_*=\frac{6c_{\Gamma^*}}{\alpha}\Big).
\]
Suppose $n>\tau_1\tilde c\,d\,p\max\{c_w^2,1/4\}$ for some $\tau_1>1$, and that on $\mathcal E_w$
\[
\lambda\,\underline w>\frac{3\kappa}{\alpha}\max\{c_{M^*},1\}\,\varepsilon .
\]
For the implemented weights ($\hat w\ge1$) it suffices that $\lambda>\frac{3\kappa}{\alpha}\max\{c_{M^*},1\}\varepsilon$; this is DDK's condition with $2-\alpha$ replaced by $\kappa$. Then, with probability at least $1-c_2\exp(-\tau_1p)-P(\mathcal E_w^c)$:
\begin{enumerate}[label=\upshape(\alph*)]
  \item the minimiser of \eqref{eq:est} is unique, $\hat S\subseteq S$, and
  \[
  \|\hat M-M^*\|_\infty\le\frac{c_{\Gamma^*}}{1-c_{\Gamma^*}\varepsilon}\big((c_{M^*}+1)\varepsilon+\lambda\bar w_S\big)<c_{\Gamma^*}\lambda\Big(\bar w_S+\frac\alpha\kappa\underline w\Big);
  \]
  \item if $\beta_{\min}$ exceeds the first bound in (a), then $\hat S=S$ with the signs of $M^*$.
\end{enumerate}
\end{theorem}

\begin{proof}
DDK Appendix~F, steps 1--3, with $\varepsilon_1=\varepsilon_2=\varepsilon$.
\begin{itemize}
  \item \textbf{Steps 1--2 are unchanged.} $\varepsilon/(\sqrt d(4+8c_{\Sigma^*}))$ and $\varepsilon/(2c_C)$ lie in the admissible range of Corollary~E.5. The upper end uses only $\max\{c_w^2,1/4\}\ge1/4$, the lower end only $\tilde c$. So with probability at least $1-c_2\exp(-\tau_1p)$, $\opnorm{(\Delta_\Gamma)_{\cdot S}}_\infty<\varepsilon$ and $\|\Delta_g\|_\infty<\varepsilon$. Both are bounds on the same random variable $\opnorm{\Delta_\Sigma}_2$, so no union bound is needed.
  \item \textbf{Step 3 replaces $c_*$ by $c_w$.} $\varepsilon<1/\max\{c_w,1/2\}\le1/c_w=\alpha/(3c_{\Gamma^*}(\kappa\bar r+\alpha))\le\alpha/(3c_{\Gamma^*}(\kappa r+\alpha))$ on $\mathcal E_w$, which is (W3).
  \item \textbf{Conclude.} The $\lambda$-condition is (W4), since $\max\{c_{M^*}\varepsilon,\varepsilon\}=\max\{c_{M^*},1\}\varepsilon$. Apply Theorem~\ref{thm:W1} on $\mathcal E_w\cap\{\text{concentration}\}$.\qedhere
\end{itemize}
\end{proof}

\emph{Reading.} With $\hat w\equiv1$, $\kappa=2-\alpha$ and $\bar r=1$, Theorem~\ref{thm:W2} is DDK's Theorem~2. With a good pilot, $\bar r\to0$ and one may take $\alpha\to1$, so $c_w\to3c_{\Gamma^*}$. The sample size then no longer depends on an irrepresentability gap: it becomes $n\gtrsim\tilde c\,dp\cdot9c_{\Gamma^*}^2$, against DDK's $\tilde c\,dp\cdot36c_{\Gamma^*}^2/\alpha^2$ with $\alpha$ the (often non-existent) lasso gap. The order $\Omega(dp)$ in $(d,p)$ is unchanged.
% =========================================================================
\section{Corollaries, necessity, and the dense-end pilot}\label{sec:corollaries}

\begin{corollary}[A pilot close to the truth]\label{cor:W3}
Let $\delta < \beta_{\min}/(1+\rho)$ and put
\[
  \bar r := \frac{\delta}{\beta_{\min}-\delta},
  \qquad
  \alpha := 1-\rho\,\bar r .
\]
Then $\alpha \in (0,1]$, and on the event
$\mathcal E_\delta := \{\|\hat M^0-M^*\|_\infty \le \delta\}$:
\begin{itemize}
  \item[(i)] $\hat M^0_k \neq 0$ for every $k \in S$; hence $\bar w_S < \infty$,
        and every infinite weight lies in $S^c$;
  \item[(ii)] $r \le \bar r$;
  \item[(iii)] $\rho_w \le \rho\, r \le 1-\alpha$, i.e.\ (W2) holds with this $\alpha$.
\end{itemize}
Consequently, \textbf{whatever the size of DDK's constant $\rho$}:
\begin{itemize}
  \item[(a)] on $\mathcal E_\delta$, the conclusions of Theorem~\ref{thm:W1} hold with this
        $\alpha$ whenever (W1), (W3) and (W4) hold; since $r \le \bar r$, (W3) is implied by
        $\varepsilon_1 \le \alpha / \big(3c_{\Gamma^*}(\kappa\bar r + \alpha)\big)$;
  \item[(b)] Theorem~\ref{thm:W2} applies with $\mathcal E_w := \mathcal E_\delta$ and with this
        $\alpha$ and $\bar r$: under its sample-size and $\lambda$-conditions, its conclusions hold
        with probability at least
        $1 - c_2\exp(-\tau_1 p) - P\big(\|\hat M^0-M^*\|_\infty > \delta\big)$.
\end{itemize}
\end{corollary}

\begin{proof}
Throughout, work on $\mathcal E_\delta$, so that
$|\hat M^0_j - M^*_j| \le \|\hat M^0 - M^*\|_\infty \le \delta$ for every entry $j$.

\emph{Step 0: $\alpha \in (0,1]$.}
Since $\rho \ge 0$, we have $\delta < \beta_{\min}/(1+\rho) \le \beta_{\min}$, so
$\beta_{\min} - \delta > 0$, $\bar r \ge 0$ is well defined, and $\alpha \le 1$. Moreover,
multiplying by $\beta_{\min} - \delta > 0$,
\[
  \alpha > 0
  \iff \rho\,\frac{\delta}{\beta_{\min}-\delta} < 1
  \iff \rho\,\delta < \beta_{\min} - \delta
  \iff (1+\rho)\,\delta < \beta_{\min},
\]
and the last inequality is the assumption.

\emph{Step 1: location of the pilot entries; proof of (i).}
Off the support: for $j \in S^c$ we have $M^*_j = 0$, hence
\begin{equation}\label{eq:W3-off}
  |\hat M^0_j| \;=\; |\hat M^0_j - M^*_j| \;\le\; \delta .
\end{equation}
On the support: for $k \in S$, the triangle inequality
$|M^*_k| \le |M^*_k - \hat M^0_k| + |\hat M^0_k|$ gives the reverse triangle inequality
\begin{equation}\label{eq:W3-on}
  |\hat M^0_k| \;\ge\; |M^*_k| - |\hat M^0_k - M^*_k| \;\ge\; \beta_{\min} - \delta \;>\; 0,
\end{equation}
where we used $|M^*_k| \ge \beta_{\min}$ (definition of $\beta_{\min}$ as the minimum over $S$),
$|\hat M^0_k - M^*_k| \le \delta$, and $\delta < \beta_{\min}$.
Hence $\hat M^0_k \neq 0$ for every $k \in S$, so every weight on $S$ falls into the first case
of the weight definition, $\hat w_k = \hat\mu / |\hat M^0_k|$, with $\hat\mu < \infty$ and a
strictly positive denominator. Thus each $\hat w_k$, $k \in S$, is finite, and
$\bar w_S = \max_{k \in S} \hat w_k$, a maximum of finitely many finite numbers, is finite.
An infinite weight occurs only at a pilot zero, which by \eqref{eq:W3-on} is impossible on $S$.
Hence the constraint $m_j = 0$ imposed by $\hat w_j = \infty$ only fixes entries with
$M^*_j = 0$; no entry of the true support is excluded from the estimator.

\emph{Step 2: proof of (ii).}
By \eqref{eq:W3-on}, $\hat\mu \ge |\hat M^0_k| \ge \beta_{\min} - \delta > 0$ for any $k \in S$.
With the convention $\hat\mu/0 := \infty$, which matches the definition of the excluded weights,
$\hat w_j = \hat\mu / |\hat M^0_j|$ for every entry, and $t \mapsto \hat\mu/t$ is decreasing on
$[0,\infty)$. Hence the largest weight on $S$ belongs to the smallest pilot entry on $S$, and
the smallest weight on $S^c$ belongs to the largest pilot entry on $S^c$:
\[
  \bar w_S = \frac{\hat\mu}{\min_{k\in S}|\hat M^0_k|},
  \qquad
  \underline w = \frac{\hat\mu}{\max_{j\in S^c}|\hat M^0_j|}.
\]
(If $\hat M^0_{S^c} = 0$, then $\underline w = \infty$ and $r = 0$, and (ii) is trivial.)
Therefore
\[
  r = \frac{\bar w_S}{\underline w}
    = \frac{\max_{j\in S^c}|\hat M^0_j|}{\min_{k\in S}|\hat M^0_k|}
    \le \frac{\delta}{\beta_{\min} - \delta} = \bar r,
\]
because by \eqref{eq:W3-off} the numerator is at most $\delta$, by \eqref{eq:W3-on} the
denominator is at least $\beta_{\min} - \delta > 0$, and such a ratio increases when its
numerator increases or its denominator decreases.

\emph{Step 3: proof of (iii).}
Let $j \in S^c$. If $\hat w_j = \infty$, the $j$-th term in the definition (1.2) of $\rho_w$
is $0$. Otherwise, using $\hat w_j \ge \underline w$, $\hat w_k \le \bar w_S$ for $k \in S$,
and $\sum_{k\in S}|P_{jk}| \le |||P|||_\infty = \rho$ (maximum absolute row sum),
\[
  \frac{1}{\hat w_j}\sum_{k\in S}|P_{jk}|\,\hat w_k
  \;\le\; \frac{\bar w_S}{\underline w}\sum_{k\in S}|P_{jk}|
  \;\le\; r\,\rho .
\]
Taking the maximum over $j \in S^c$ gives $\rho_w \le \rho\, r$. By (ii) and the definition
of $\alpha$, $\rho_w \le \rho\, r \le \rho\,\bar r = 1 - \alpha$, and $\alpha \in (0,1]$ by
Step 0. This is (W2).

\emph{Step 4: consequences.}
(a) By (i), the standing assumption $\bar w_S < \infty$ of Theorem~\ref{thm:W1} holds, and by
(iii) so does (W2). The remaining hypotheses (W1), (W3), (W4) are assumed. The right-hand side
of (W3) is decreasing in $r$ and $r \le \bar r$ by (ii), so
$\varepsilon_1 \le \alpha/\big(3c_{\Gamma^*}(\kappa\bar r + \alpha)\big)$ implies (W3).

(b) Theorem~\ref{thm:W2} requires an event on which $\bar w_S < \infty$,
$\rho_w \le 1 - \alpha$ and $r \le \bar r$, for deterministic $\alpha$ and $\bar r$.
By (i)--(iii), $\mathcal E_\delta$ is such an event, and $\alpha$, $\bar r$ depend only on
$\delta$, $\beta_{\min}$ and $\rho$. Theorem~\ref{thm:W2} therefore gives its conclusions with
probability at least $1 - c_2\exp(-\tau_1 p) - P(\mathcal E_\delta^c)$.
\end{proof}

\begin{remark}\label{rem:W3}
\begin{itemize}
  \item Since the diagonal is penalised, $\beta_{\min} = \min_{(j,k)\in S}|M^*_{jk}|$ runs over
        all of $S$, diagonal included. The diagonal entries are typically the largest in
        magnitude, so in practice $\beta_{\min}$ is the weakest true edge, and the tolerance
        $\beta_{\min}/(1+\rho)$ is set by that edge.
  \item DDK need $\rho \le 1-\alpha < 1$, a property of $(\Sigma^*, S)$ that no amount of data
        changes. Here $\rho$ only determines how accurate the pilot must be: for every $\rho$
        the tolerance $\beta_{\min}/(1+\rho)$ is positive. The constant $\rho$ does not
        disappear, however. It enters $\kappa \ge 1+\rho$ in (W3)--(W4) and
        $\alpha = 1 - \rho\bar r$, and as $\delta \uparrow \beta_{\min}/(1+\rho)$ we have
        $\alpha \downarrow 0$, so (W3), (W4) and the sample size in Theorem~\ref{thm:W2}
        become arbitrarily demanding.
  \item For a prescribed $\alpha \in (0,1]$, the condition $\rho\bar r \le 1-\alpha$ is
        equivalent to $\delta \le (1-\alpha)\beta_{\min}/(\rho + 1 - \alpha)$. The $\alpha$ of
        Corollary~\ref{cor:W3} is the largest admissible one, and also the best choice, since
        (W3) and (W4) become weaker requirements as $\alpha$ grows.
\end{itemize}
\end{remark}


\begin{corollary}[Sign-weighted version]\label{cor:W4}
Assume (W1), (W3), (W4) and $\beta_{\min}>b_w$. Replace (W2) by the weaker condition
\[
\rho_w^{\mathrm{sgn}}:=\max_{j\in S^c}\frac{1}{\hat w_j}\Big|\sum_{k\in S}P_{jk}\,\hat w_k\sign(M^*_k)\Big|\le1-\alpha ,
\]
where the sum includes the diagonal, with $\sign(M^*_{ii})=-1$ in the thesis DGP. Then the conclusions of Theorem~\ref{thm:W1} hold.
\end{corollary}

\begin{proof}
\eqref{eq:errbound} does not use (W2). With $\beta_{\min}>b_w$ it gives $\sign(\tilde m_S)=\sign(m^*_S)$, hence $\tilde z_S=\sign(m^*_S)$. The deterministic term of \eqref{eq:dual} then equals the expression inside $\rho_w^{\mathrm{sgn}}$; the rest of the proof is unchanged.
\end{proof}

This is the weighted form of DDK's weak condition (G.1). It is the quantity to compute per dataset; \code{gclmth.adaptive_conditions} returns it for an unpenalised diagonal and needs the diagonal signs added for this version.

\begin{proposition}[Necessity; the analogue of DDK Prop.~G.2]\label{prop:W5}
Let $\lambda>0$ be fixed and let the weights converge in probability, $\hat w\to w$, with $0<w_j<\infty$ on $S^c$ and $w<\infty$ on $S$. Assume $\beta_{\min}>c_{\Gamma^*}\lambda\max_{k\in S}w_k$ and
\[
\max_{j\in S^c}\frac{1}{w_j}\Big|\sum_{k\in S}P_{jk}\,w_k\sign(M^*_k)\Big|>1 .
\]
Then $P(\hat S=S)\to0$.
\end{proposition}

\begin{proof}
As in G.2. On $\{\hat S=S\}$, $\hat M$ is the restricted solution \eqref{eq:restricted}, and by \eqref{eq:errbound} and $\Delta_\Gamma,\Delta_g\to0$ in probability (Lemma~E.6, Corollary~E.5) it has the signs of $M^*$ with probability tending to one. Its dual vector is \eqref{eq:dual}. The noise part tends to $0$, and the deterministic part tends to the displayed quantity, which exceeds $1$ in some row. That violates \eqref{eq:kkt}, so $P(\hat S=S)\to0$.
\end{proof}

\begin{remark}[The dense-end pilot when the diagonal is penalised]\label{rem:W6}\leavevmode
\begin{itemize}
  \item As $\lambda\to0$, and in the population, the pilot tends to the minimum-$\|M\|_1$ exact solution (all entries) of $M\Sigma^*+\Sigma^*M^\top+C=0$ (Tibshirani 2013). Penalising the diagonal changes this basis-pursuit problem, but only in one direction.
  \item \emph{Trace identity.} Every exact solution has $\tr(M)=-\tfrac12\tr(C\Sigma^{*-1})$. To see this, multiply the equation by $\Sigma^{*-1}$ from the right and take traces, using $\tr(\Sigma^*M^\top\Sigma^{*-1})=\tr(M)$. Hence $\sum_i|M_{ii}|\ge\tfrac12\tr(C\Sigma^{*-1})$, with equality if and only if all $M_{ii}\le0$.
  \item \emph{Consequence.} Suppose $M^*$ has a non-positive diagonal (every DGP of the thesis does) and is the unique minimiser of $\|M_{\mathrm{off}}\|_1$ over exact solutions; this is the BP condition checked so far. Then for every exact $M\ne M^*$,
  \[
  \|M\|_1\ge\|M_{\mathrm{off}}\|_1+\tfrac12\tr(C\Sigma^{*-1})>\|M^*_{\mathrm{off}}\|_1+\tfrac12\tr(C\Sigma^{*-1})=\|M^*\|_1 .
  \]
  So the off-diagonal BP condition implies the penalised-diagonal one. In the population, penalising the diagonal can only help the dense-end pilot, and the BP counts computed so far are lower bounds.
  \item For an exact check of the penalised version, the dual certificate requires $A_S^\top h=\sign(m^*_S)$ on all of $S$, diagonal included, instead of $A_D^\top h=0$ on the diagonal. This is a one-line change in \code{bp_certificate}.
  \item When the BP condition holds, $\hat M^0\to M^*$, so $E\supseteq S^c$ in the limit, $r\to0$, $\rho_w\to0$, and Corollary~\ref{cor:W3} applies.
\end{itemize}
\end{remark}

% =========================================================================
\section{What changed relative to DDK, line by line}\label{sec:table}

{\small
\renewcommand{\arraystretch}{1.6}
\noindent\begin{tabularx}{\textwidth}{>{\raggedright\arraybackslash}p{0.27\textwidth}>{\raggedright\arraybackslash}p{0.36\textwidth}>{\raggedright\arraybackslash}X}
\toprule
DDK (lasso, $\hat w\equiv1$) & here (weights $\hat w_j=\hat\mu/|\hat M^0_j|\ge1$) & why\\
\midrule
(D.1) $\rho<1-\alpha$
& (W2) $\rho_w\le1-\alpha$, with $\rho_w\le\rho r$
& the deterministic part of the dual vector \eqref{eq:det} is $\hat w_j^{-1}[P(\hat w\circ\tilde z)_S]_j$\\
$\varepsilon_1\le\alpha/(6c_{\Gamma^*})$
& (W3) $\varepsilon_1\le\alpha/(3c_{\Gamma^*}(\kappa r+\alpha))$
& $G_2$ carries $\kappa r$: bias $\lambda\bar w_S$ on $S$ against kink $\lambda\underline w$ on $S^c$\\
$\lambda>\frac{3(2-\alpha)}{\alpha}\max\{c_{M^*}\varepsilon_1,\varepsilon_2\}$
& (W4) $\lambda\underline w>\frac{3\kappa}{\alpha}\max\{c_{M^*}\varepsilon_1,\varepsilon_2\}$; implied by $\lambda>\frac{3\kappa}{\alpha}\max\{\cdot\}$ since $\hat w\ge1$
& noise is compared with the smallest effective kink; $2-\alpha\to\kappa\ge1+\rho$ because $\rho$ may exceed $1$\\
(D.4) $\|\tilde\Delta_M\|_\infty<\frac{2c_{\Gamma^*}\lambda}{2-\alpha}$
& \eqref{eq:errbound}--\eqref{eq:bwlambda} $\le b_w<c_{\Gamma^*}\lambda(\bar w_S+\frac\alpha\kappa\underline w)$
& bias $\lambda\to\lambda\bar w_S$\\
(D.5)
& \eqref{eq:dual}
& $\sign(\cdot)\to\hat w\circ\tilde z$, row $j$ divided by $\hat w_j$\\
beta-min $>\frac{2c_{\Gamma^*}\lambda}{2-\alpha}$
& $\beta_{\min}>b_w$ (all of $S$, diagonal included)
& same role\\
Thm~2: $c_*=6c_{\Gamma^*}/\alpha$
& $c_w=3c_{\Gamma^*}(\kappa\bar r+\alpha)/\alpha$
& from (W3)\\
Prop.~G.2
& Prop.~\ref{prop:W5}
& dual vector with limiting weights\\
\bottomrule
\end{tabularx}}

\medskip
Everything specific to GCLMs enters exactly where it enters in DDK: through the invertibility of $\Gamma^*_{SS}$, the constants $c_{\Gamma^*},\rho,c_{M^*},c_{\Sigma^*},c_C$, and Appendix~E. Invertibility of $\Gamma^*_{SS}$ holds for every simple graph (Dettling et al.\ 2023).

% =========================================================================
\section{Remarks}\label{sec:remarks}

\paragraph{Numerical check} (\code{check_W1.py}, \code{gclmth}; weighted lasso with the diagonal penalised and the implemented weights).
\begin{itemize}
  \item \emph{In-star, $p=6$, $|w|=0.7$, correlation scale.} DDK's constant is $\rho=1.78$, so the lasso's condition fails. The pilot is truth $+\,U(-0.05,0.05)$, giving $\rho_w=0.066$ and $r=0.052$. At $n=10^6$ and $10^7$, (W3) holds. With $\lambda$ at (W4) plus 1\,\%, the weighted lasso returned exactly $S$, with $\|\hat M-M^*\|_\infty=0.13$ and $0.057$, below $b_w=0.79$ and $0.31$, which are below the right-hand side of \eqref{eq:bw}.
  \item \emph{Unweighted 3-path with ordered diagonal, raw scale ($\rho=0.81$, DDK's case), $\kappa=2-\alpha$.} The right-hand side of \eqref{eq:bw} equals DDK's $2c_{\Gamma^*}\lambda/(2-\alpha)$ to machine precision. $\hat S\subseteq S$ held, and part (b) was not reached because $\lambda$ is large relative to $\beta_{\min}=0.2$.
\end{itemize}
The required $n$ is large because DDK's $\varepsilon_1$ (row sums of $\Delta_\Gamma$ over the support columns) is a conservative quantity. The theorems are sufficient conditions, as in DDK.

\paragraph{Unpenalised diagonal} (the default runs, waves 1--9). Set $\hat w=0$ on the diagonal $D$, take the maximum defining $\hat\mu$ over off-diagonal entries, replace $S$ by $S\setminus D$ in $\bar w_S$, \eqref{eq:rhow}, Corollary~\ref{cor:W4} and the beta-min, and read the diagonal rows of \eqref{eq:kkt} as $(\hat\Gamma\hat m-\hat g)_D=0$. Nothing else changes; the proof is the one above.

\paragraph{Points in DDK to be aware of when citing.}
\begin{enumerate}
  \item $g$ should carry the transpose (\S\ref{sec:setting}).
  \item D.1's $\lambda$-condition is printed as $\max\{c_{M^*},\varepsilon_1,\varepsilon_2\}$; the proof uses $\max\{c_{M^*}\varepsilon_1,\varepsilon_2\}$ ($G_1$, $G_3$ on pp.~532--533), and Theorem~2 uses $c_{M^*}\varepsilon$ (implicitly $c_{M^*}\ge1$).
  \item The beta-min displays index $(j,k)\in S$ with $1\le j<k$, inherited from the symmetric setting of Lin et al.; for a drift matrix it must run over all of $S$.
  \item The block form on p.~531 writes $+\Delta_g$; with $\Delta_g=\hat g-g^*$ the sign is $-\Delta_g$. Only norms enter, so nothing depends on it.
\end{enumerate}

\end{document}
